\section{Dyskusja}\label{Dyskusja}

% Dyskusja

Celem niniejszej pracy była analiza postaw, wiedzy oraz praktyk studentów w zakresie stosowania suplementów diety. Uzyskane wyniki wskazują na niezwykle wysoką powszechność tego zjawiska w badanej grupie – aż 90,7\% respondentów zadeklarowało aktualne przyjmowanie preparatów uzupełniających dietę. Wynik ten jest znacząco wyższy niż odnotowany w badaniu Sicińskiej i wsp. (2022) na Uniwersytecie Warszawskim (41\%) czy w badaniach studentów w Chorwacji (30,5\%) \cite{diet-lifestyle, croatia-knowledge-and-use}. Bliższy jest natomiast danym z Arabii Saudyjskiej (76,6\%) czy Bejrutu (68\%) \cite{alfawaz2017prevalence, pharmacist-knowledge-and-use}. Tak wysoki odsetek w badaniu własnym wynika prawdopodobnie ze specyfiki grupy badawczej, w której dominowali studenci psychologii i dietetyki (łącznie ponad 72\%), co zazwyczaj wiąże się z większym zainteresowaniem zdrowiem, profilaktyką oraz dobrostanem psychofizycznym.

W literaturze przedmiotu płeć biologiczna jest wskazywana jako jeden z najsilniejszych predyktorów stosowania suplementów. Badania Sicińskiej (2022) czy Vidovica (2022) sugerują, że mężczyźni sięgają po suplementy znacznie rzadziej niż kobiety (o ok. 38\%) \cite{sicinska2022dietary, vidovic2022dietary}. Co ciekawe, badania własne nie potwierdziły tej zależności ($p = 0,36$). Nie wykazano statystycznie istotnych różnic w częstości suplementacji między kobietami (88,9\%) a mężczyznami (100,0\%). Może to świadczyć o zacieraniu się różnic międzypłciowych w grupach o podobnym profilu wykształcenia (studenci akademicy) i wysokich zainteresowaniach prozdrowotnych, choć interpretację tę należy prowadzić ostrożnie ze względu na dużą dysproporcję liczebną płci w badanej próbie (84,1\% kobiet vs 15,9\% mężczyzn), co mogło ograniczyć moc testu statystycznego.

Podobny brak istotności odnotowano w przypadku wieku ($p = 0,54$) oraz poziomu aktywności fizycznej ($p = 0,17$). Choć we wszystkich podgrupach większość osób stosowała suplementy, to brak związku z aktywnością fizyczną stoi w lekkiej sprzeczności z częścią literatury (np. \cite{diet-lifestyle}), gdzie intensywny styl życia i sport mocno determinowały suplementację. W badanej grupie niska frekwencja stosowania typowych odżywek sportowych, takich jak białko (24,7\%) czy kreatyna (15,4\%), koreluje z faktem, że większość badanych (33,6\%) podejmowała aktywność rzadziej niż raz w tygodniu, a dominującą część próby stanowiły kobiety, które według literatury rzadziej wybierają preparaty ukierunkowane na budowanie masy mięśniowej \cite{vidovic2022dietary, radwan2019prevalence}.

Dominującym preparatem w badanej grupie była witamina D (81,4\%), co jest w pełni uzasadnione szeroką kampanią edukacyjną w Polsce dotyczącą powszechnych niedoborów tego składnika w naszej strefie klimatycznej. Na kolejnych miejscach znalazły się kwasy omega-3 (60,8\%) oraz magnez (49,4\%). Wyniki te różnią się od danych z USA czy Jordanii, gdzie studenci najczęściej wybierają ogólne preparaty multiwitaminowe i multimineralne \cite{lieberman2015patterns, knowledge-and-use}.

Główną motywacją badanych była chęć poprawy zdrowia (70,1\%), wzmocnienia odporności (68,04\%) oraz zapobiegania niedoborom (67,01\%). Koresponduje to z trendami światowymi opisywanymi przez Sundgota-Borgen i wsp. (2022) oraz badaniami z Bejrutu, gdzie motywacje zdrowotne i profilaktyczne wysuwały się na pierwszy plan \cite{sundgot2022explanations, pharmacist-knowledge-and-use}. Warto zauważyć, że blisko połowa respondentów (47,42\%) wskazała również na powody estetyczne (poprawa wyglądu skóry, włosów i paznokci), co idealnie wpisuje się w doniesienia Lawanda i współpracowników (2023) o rosnącym znaczeniu aspektów kosmetycznych w decyzjach o suplementacji wśród młodych dorosłych \cite{lawand2023assessment}.

Głównym źródłem informacji o suplementach diety dla badanych studentów pozostał internet (69,07\%), co jest zjawiskiem powszechnym i zgodnym z badaniami globalnymi, gdzie media cyfrowe osiągają wskaźniki rzędu 38\%–40\% \cite{kobayashi2017prevalence, chiba2020educational}. Niepokojący w literaturze jest fakt, że rzadko towarzyszą temu konsultacje ze specjalistami \cite{croatia-knowledge-and-use}. W badaniach własnych odsetek konsultacji z lekarzem (40,21\%) i dietetykiem (38,14\%) okazał się relatywnie wysoki na tle innych krajów, co ponownie można przypisać obecności studentów dietetyki w próbie.

Zauważalny jest jednak dualizm postaw: lekarz jest uznawany przez studentów za najbardziej wiarygodne źródło informacji (69,07\%), a naukowe portale internetowe zajmują drugie miejsce (40,21\%). Mimo tak wysokiego zaufania deklaratywnego, w praktyce aż 47,42\% respondentów wskazało, że w ogóle nie konsultuje swojej suplementacji ze specjalistą, podejmując decyzje autonomicznie. Zjawisko to potwierdza teorię Steele i Senekal (2005) o silnym wpływie marketingu bezpośredniego i subiektywnego poczucia bezpieczeństwa stosowania tych produktów \cite{steele2005dietary}.

Kluczowym i najbardziej intrygującym wnioskiem z przeprowadzonej analizy statystycznej jest brak przełożenia wiedzy na realne praktyki bezpiecznej suplementacji. Z jednej strony, studenci wykazali się dobrą zdolnością do samokrytyki – ich samoocena wiedzy umiarkowanie, ale istotnie korelowała z rzeczywistym wynikiem testowym (dla ogólnej wiedzy $\rho = 0,47$, dla interakcji $\rho = 0,54$, przy $p < 0,001$). Średni obiektywny poziom wiedzy wyniósł 3,76 na 6 punktów, a blisko 40\% badanych osiągnęło wysoki próg punktowy. Wynik ten jest lepszy niż w badaniach Homana (2018) czy Elsahoryi (2023), gdzie studenci odpowiadali poprawnie na mniej niż połowę pytań \cite{homan2018dietary, elsahoryi2023prevalence}.

Z drugiej strony, analiza wykazała brak istotnej korelacji między indeksem wiedzy a indeksem prawidłowych praktyk suplementacyjnych ($\rho = 0,15$; $p = 0,11$). Co więcej:

\begin{itemize}
    \item Wyższy poziom wiedzy nie różnicował osób pod kątem wykonywania badań krwi przed rozpoczęciem suplementacji ($p = 0,19$).
    \item Wiedza nie wpływała na to, czy studenci konsultowali się ze specjalistą ($p = 0,57$).
    \item Zależność między wiedzą a przekraczaniem zalecanych dawek znalazła się na granicy istotności statystycznej ($p = 0,052$), co sugeruje jedynie słaby trend, według którego wyższa świadomość może (choć nie musi) chronić przed dawkowaniem niezgodnym z zaleceniami.
\end{itemize}

Jedynym obszarem, w którym wiedza odegrała skrajnie istotną rolę, było sprawdzanie powielania składników w przyjmowanych preparatach ($p = 0,0001$). Średnia wiedzy w grupie sprawdzającej skład wyniosła 4,9 pkt. względem 3,6 pkt. w grupie niesprawdzającej.

Ten stan rzeczy obnaża wyraźną lukę behawioralną: studenci posiadający wiedzę teoretyczną potrafią zastosować ją do czytania etykiet i unikania dublowania dawek (co jest czynnością prostą, analityczną), jednak ta sama wiedza nie motywuje ich do wdrażania głębszych procedur bezpieczeństwa, takich jak diagnostyka laboratoryjna czy profesjonalna konsultacja medyczna. Wpisuje się to w mity opisywane przez Chibę i wsp. (2020), zgodnie z którymi młode osoby, nawet dobrze wyedukowane, podświadomie traktują suplementy jako „całkowicie bezpieczną żywność naturalną”, przez co rezygnują z rygorystycznych środków ostrożności \cite{chiba2020educational}. Ogólny indeks praktyk w badaniu własnym okazał się niski u aż 41,24\% respondentów, co potwierdza, że wiedza teoretyczna nie jest wystarczającym czynnikiem gwarantującym bezpieczeństwo zdrowotne.

Głównym ograniczeniem niniejszego badania jest brak pełnej reprezentatywności akademickiej – silne skrzywienie próby w stronę kobiet oraz studentów kierunków o profilu medycznym i humanistycznym (dietetyka, psychologia) utrudnia generalizację wyników na całą populację studencką (np. kierunki techniczne czy ścisłe). Ponadto mała liczebność grupy niestosującej suplementów ($n = 10$) ograniczyła możliwości porównawcze w niektórych testach.

Mimo to, uzyskane dane wyraźnie definiują potrzeby edukacyjne młodych dorosłych. Prawie połowa studentów (50,12\%) nie wiedziała o braku konieczności przeprowadzania badań klinicznych dla suplementów diety przed wprowadzeniem ich na rynek, co potwierdza tezę Homana (2018) o nieznajomości podstaw prawnych i regulacyjnych \cite{homan2018dietary}. Projektując przyszłe interwencje edukacyjne, należy porzucić suche wykłady akademickie na rzecz form preferowanych przez samych studentów – ponad połowa (54,64\%) wskazała krótkie filmy edukacyjne, a duża część artykuły popularnonaukowe (46,39\%) i warsztaty praktyczne (40,21\%). Edukacja ta powinna kłaść nacisk nie tyle na biochemiczną teorię witamin, ile na kształtowanie prawidłowych nawyków behawioralnych (konieczność badań krwi, ryzyko samowolnej suplementacji).

